\documentclass[11pt,a4paper]{article}
\usepackage[margin=25mm]{geometry}
\usepackage{amsmath,amssymb,amsthm,hyperref}
\newtheorem{theorem}{Theorem}
\newtheorem{lemma}[theorem]{Lemma}
\newtheorem{corollary}[theorem]{Corollary}
\newcommand{\Sym}{\operatorname{Sym}}
\newcommand{\Core}{\operatorname{Core}}
\title{Reducing permutation degree through a commuting action}
\author{Human Compression project: working common argument}
\date{}
\begin{document}\maketitle

A free transitive group action supplies coordinates without selecting an
identity point intrinsically: it is a torsor. Choosing one point in each
orbit of a semiregular permutation group makes the action of its centralizer
particularly explicit. A quotient of these coordinates can reduce the
permutation degree while preserving the entire commuting group, including
its marked subgroup intersections. The quotient need not preserve the
original geometry; the assertion concerns a faithful representation of
the same abstract group.

\begin{lemma}\label{free}
Let $H\leq\Sym(\Omega)$ be transitive and let
$D\leq C_{\Sym(\Omega)}(H)$. Then $D$ is semiregular.
\end{lemma}
\begin{proof}
If $d\in D$ fixes $x$, then it fixes every $h(x)$, because $dh=hd$.
Transitivity implies that $d$ fixes $\Omega$ pointwise. Since $D$ is a
permutation group, $d=1$.
\end{proof}

\begin{theorem}\label{reduce}
Let a finite group $D$ act semiregularly on a nonempty finite set $\Omega$,
and let $b$ be its number of orbits. Suppose $K\leq D$ has trivial core,
\[
                \Core_D(K)=\bigcap_{d\in D}d^{-1}Kd=1.
\]
There is a faithful permutation representation of
$C_{\Sym(\Omega)}(D)$ on $b[D:K]$ points. Every transitive subgroup
$H$ of this centralizer remains transitive in the new representation.
If $b>1$ and $[D:K]>1$, the new action preserves a nontrivial system of
$b$ blocks of size $[D:K]$.
\end{theorem}
\begin{proof}
Choose one point $x_j$ in each $D$-orbit, indexed by a set $B$ of size $b$.
Semiregularity identifies $\Omega$ with $D\times B$ by
$(d,j)\mapsto d(x_j)$. The action of $D$ in these coordinates is left
multiplication in the first coordinate.

A permutation $h$ commutes with $D$ if and only if there are
$\sigma\in\Sym(B)$ and elements $a_j\in D$ such that
\[
                 h(d,j)=(d a_j,\sigma(j)).
\]
Indeed, its values at the $x_j$ determine all other values by commutation.
Conversely every displayed map is a bijection and commutes with left
multiplication. Composing two such maps again gives a map of the same form.
This is the usual wreath description of the centralizer, with the right
regular action in each coordinate.

Now let $\Delta=K\backslash D$ be the right $D$-set of right cosets, and
define the action on $\Delta\times B$ by
\[
             \widehat h(Kd,j)=(Kd a_j,\sigma(j)).
\]
It is well defined, and the same composition formula shows that
$h\mapsto\widehat h$ is a homomorphism. If $\widehat h$ fixes every
point, then $\sigma=1$, and every $a_j$ acts trivially on $K\backslash D$.
For any $a\in D$, this last condition is equivalent to
$d a d^{-1}\in K$ for every $d\in D$, hence to $a\in\Core_D(K)$.
Thus all $a_j=1$, proving faithfulness.

The surjection $(d,j)\mapsto(Kd,j)$ is equivariant for every commuting
$h$. Therefore a transitive $H$ remains transitive on its image. Each
$\Delta\times\{j\}$ maps onto $\Delta\times\{\sigma(j)\}$, proving
the block assertion. All changes of the chosen orbit representatives
give conjugate descriptions of this construction.
\end{proof}

\begin{corollary}\label{marked}
In Theorem~\ref{reduce}, replace a marked subgroup
$(H;(s_i)_{i\in I})$ by $(\widehat H;(\widehat s_i)_{i\in I})$.
Generator orders, product orders, irredundance and every marked subgroup
intersection identity are unchanged. In particular string C-group
representations retain their rank and intersection property.
\end{corollary}
\begin{proof}
The representation restricts to an isomorphism from $H$ to its image.
It sends $\langle s_i:i\in J\rangle$ onto
$\langle\widehat s_i:i\in J\rangle$ for every $J\subseteq I$.
Injectivity implies that images commute with intersections of subgroups.
These facts give each assertion, including all pairs of subsets in the
C-group intersection property.
\end{proof}

\begin{corollary}\label{dihedral}
Suppose a transitive group $H\leq\Sym(\Omega)$ centralizes a dihedral
group $D$ of order $2d$, where $d\geq3$. Then $H$ has a faithful transitive
permutation representation of degree $|\Omega|/2$. If $H$ is a string
C-group of rank $s$ and $|\Omega|/2=m\geq7$, the transitive rank bound
for proper subgroups of symmetric and alternating groups gives
\[
                         s\leq m/2+1.
\]
\end{corollary}
\begin{proof}
Lemma~\ref{free} makes $D$ semiregular. Take $K$ to be the stabilizer of
a vertex in the usual action of $D$ on a regular $d$-gon. This action
is faithful for $d\geq3$, so $K$ is core-free and $[D:K]=d$.
Theorem~\ref{reduce} gives degree $bd=|\Omega|/2$; Corollary~\ref{marked}
preserves the C-group representation.

If $b>1$, the new action is imprimitive, so its image is neither $S_m$
nor $A_m$. If $b=1$, the centralizer of the regular $D$-action has
order $|D|$. A subgroup of it transitive on $|D|$ points must be the
entire centralizer. Thus $H\cong D$ has order $2m$, which for $m\geq7$
is smaller than $|A_m|$. In both cases the stated transitive rank theorem
applies. That theorem is a retained premise here, with its complete proof
in the accepted common theory for Topic B; it is not reproved by this note.
\end{proof}

For example, let $\Gamma$ be a string C-group of rank $r$ on $2m$ points,
and let a subgroup obtained by deleting two marked generators be
transitive and commute with two noncommuting involutions. These involutions
generate a dihedral group of order at least six. If $m\geq7$, the subgroup
has rank $r-2$, so the corollary gives
\[
                     r\leq m/2+3<m.
\]
Consequently this configuration cannot occur when $r\geq m$.
The conclusion requires actual transitivity of the subgroup; an
intransitive action does not satisfy the hypothesis.

\section*{Attribution and scope}
The centralizer of a regular action and its wreath description are standard
permutation group arguments. Cameron, Fernandes, Leemans and Mixer use
the dihedral reduction twice in Section~3.2 of \emph{Highest rank of a
polytope for $A_n$}, arXiv:1605.09173v1 (2016), pages~6--7. The general
core-free formulation above makes the reusable hypothesis and faithful
quotient map explicit. No historical priority is claimed for it.

This working note does not recover that paper's main alternating-group
rank theorem. In particular its intransitive cases, two-fracture arguments,
finite exceptions and lower bounds remain separate obligations. No whole
original, compression factor, conjecture answer or Lean formalization is
claimed. A later application must acquire every common theorem it uses;
the smaller permutation action is not a discount in manuscript length.
\end{document}
